\documentclass[11pt,a4paper]{article}
\usepackage{curso}
\begin{document}
\portada{Lección}{1.4}{El demonio de Maxwell y la máquina de Szilard}

\noindent\textbf{Unidad 1.}\quad\textbf{Objetivos:} plantear la paradoja; calcular el trabajo de Szilard; explicar la
resolución de Bennett. \textbf{Prerrequisitos:} 1.1 y 1.3.

\section{La paradoja (Maxwell, 1867)}
Un «demonio» abre una puerta entre dos cámaras solo para dejar pasar moléculas rápidas de izquierda a derecha y lentas
de derecha a izquierda. Sin trabajo aparente crea una diferencia de temperatura, que podría accionar una máquina
térmica: un móvil perpetuo de segunda especie. La animación es una simulación de billar 2D \evS{} con esa regla de
puerta: $\langle v^2\rangle$ crece a la derecha y decrece a la izquierda.

\section{La máquina de Szilard (1929)}
Una molécula en una caja de volumen $V$ en contacto con un baño a $T$:
\begin{enumerate}[nosep]
\item Insertar una pared en el medio (sin trabajo, idealmente).
\item Medir el lado de la molécula: 1 bit en la memoria del demonio.
\item Usar la pared como pistón del lado correcto; expansión isotérmica lenta de $V/2$ a $V$.
\item Retirar el pistón: la molécula ocupa toda la caja; el ciclo puede repetirse.
\end{enumerate}
Con $pV=\kB T$:
\begin{equation}W_{\text{ext}}=\int_{V/2}^{V}\frac{\kB T}{V'}\dd V'=\kB T\ln2,\qquad\Delta U=0\Rightarrow Q_{\text{baño}\to\text{gas}}=W_{\text{ext}}.\end{equation}
Un bit de información se convierte en $\kB T\ln2$ de trabajo extraído del baño.

\section{Resolución: Landauer y Bennett}
Durante décadas (Szilard, Brillouin) se atribuyó el costo a la \emph{medición}. Bennett (1982) mostró que medir puede
hacerse de forma lógicamente y termodinámicamente reversible (una copia sobre un registro en estado conocido). El costo
inevitable aparece al \emph{reiniciar} la memoria del demonio para repetir el ciclo: por Landauer (lección 1.5),
$W_{\text{borrado}}\ge\kB T\ln2$. Balance por ciclo: $W_{\text{neto}}\le\kB T\ln2-\kB T\ln2=0$ \cite{bennett}.

\begin{center}
\begin{tabular}{lcccc}\toprule
Paso & gas & memoria & baño & total\\\midrule
medir & $-1$ & $+1$ & $0$ & $0$\\
expandir & $+1$ & $0$ & $-1$ & $0$\\
borrar & $0$ & $-1$ & $\ge+1$ & $\ge0$\\\midrule
ciclo & $0$ & $0$ & $\ge0$ & $\ge0$\\\bottomrule
\end{tabular}
\\[0.3em]{\small Cambios de entropía en unidades de $\kB\ln2$. En «medir», la incertidumbre del gas se convierte en
correlación con la memoria (información mutua).}
\end{center}

\section{Medición imperfecta}
Si la medición $M$ acierta con probabilidad $1-\epsilon$, el trabajo máximo extraíble por ciclo es
\begin{equation}W_{\max}=\kB T\,I(X;M)=\kB T\ln2\,[1-h(\epsilon)].\end{equation}
\begin{ejemplo}
$T=300$ K, $\epsilon=0{,}1$: $I=1-0{,}469=0{,}531$ bits, $W_{\max}\approx0{,}531\cdot2{,}87\times10^{-21}\approx1{,}5\times10^{-21}$ J.
Con $\epsilon=0{,}2$: $I=0{,}278$ bits, $W_{\max}\approx8{,}0\times10^{-22}$ J.
\end{ejemplo}

\begin{aplicacion}
Motores de información experimentales (coloides con retroalimentación, electrones individuales). Para el cómputo:
medir y controlar puede ser barato; reiniciar memorias no. Laboratorio L2 (reto): simular una máquina de Szilard y cerrar el ciclo.
\end{aplicacion}

\begin{error}
\begin{itemize}[nosep]
\item El costo inevitable está en borrar, no en medir.
\item El demonio no viola la segunda ley si se incluye su memoria.
\item La expansión debe ser cuasiestática para extraer $\kB T\ln2$.
\item Con medición ruidosa, $W_{\max}=\kB T\,I(X;M)$.
\end{itemize}
\end{error}

\section{Resumen}
Información $\to$ trabajo ($\kB T\ln2$ por bit); borrar la memoria cuesta lo mismo; la información es un recurso termodinámico.

\begin{thebibliography}{9}\small
\bibitem[9]{bennett} C. H. Bennett, Notes on Landauer's principle, reversible computation and Maxwell's demon, 2003. \url{https://arxiv.org/abs/physics/0210005}
\bibitem[10]{wolpert} D. H. Wolpert, The stochastic thermodynamics of computation, 2019.
\bibitem[2]{peliti} L. Peliti y S. Pigolotti, \emph{Stochastic Thermodynamics: An Introduction}, 2021.
\end{thebibliography}
\end{document}
